\section{Realizing a good selected point on its prescribed division lattice}

% BEGIN S77C1
\paragraph{S77C1. Moving the rows at the actual division lattice.}
Let $w$ be a column over a base, frame and twist datum $(B,f,T)$.
At a level $Z$ let $O$ be its staged object, $W_t$ the prescribed
set of places, and $t$ a row packet for $O.p_0$ and $W_t$. For each
word $u$ choose a submodule $N(u)$ of the source division module,
and choose a real height $h(u)$. Write $Z'=\mathrm{setDivZ}(Z,N,h)$.
It replaces the division data in $Z.V$, with
$w.D(Z')=\mathrm{setDiv}(w.D(Z),N,h)$. For the transported object
$O'=O.\mathrm{tr}(N,h)$ there is a row packet $t'$ at $O'.p_0$
and $W_t$ having the same rows in the sense of $\mathrm{RowsSame}$.
Its geometric component is the specified transport of $t.r.z.\mathrm{geo}$,
and its spectral row package is the upper-stage transport of the old
one. The object ties and the row-extension witness transfer to this $t'$.

Indeed apply the division change to the original geometric record
$z=t.r.z$. Apply the same substitution to its ordinary reading $Y$
and lift reading $L$ in that order. The latter now belongs to the
former. In the substituted coordinates the tuple and its twist,
the characters of $P_Y$ and $P_W$, the representation, the two eigenperiod
embeddings, $e_E,j_A,j_{A,0},e_{M_t},e_W,j_W$ and the two
Riemann--Roch coordinate equivalences agree with their originals
after identifying their types. For the last two use the
Riemann--Roch equivalences at the transported parts of $z.R$.
These are the individual fields of $\mathrm{ReadSame}$. Its row
constructor gives $t'$ and the object and row-extension implications.
Unfolding the moved record identifies its geometric component with
$\mathrm{geoHalf}$ at $Z'$ and its row package with the one moved
along $(Z,Z.S,N,h)$. All these assertions use the same $t'$.
% END S77C1

% BEGIN S77C2
\paragraph{S77C2. The position conditions survive the move.}
Fix the source local Langlands input $h_L$, depth parameter $d$
and geometric condition $\Gamma$. Write
\[
 p_0=\mathrm{prmR}\bigl(\mathrm{NEo}(h_L,d),
       \mathrm{SExt.none},\mathrm{geoCl}(\Gamma)\bigr).
\]
Given a witness $x$ at $(w,Z,O)$ for $p_0.s$, use S77C1 on its
rows for arbitrary $N,h$. The column and character data at the
moved object are arbitrary inputs of the types prescribed by
$\mathrm{PositionMv}$. We construct its four required witnesses.

The row component follows by applying the transported row-extension
implication to $x.rx$. For the character choice use the upper-stage
copy of the old choice and theta data. At $S'=Z.S$ they identify
with the choice and theta data of $O'$. The old and new tuples are
equal because their $\xi$ and $\psi$ agree. In the moved local
record the map of selected primes agrees with the old one, so its
range is still the range of $w.I.\mathrm{sel}$. The moved row is
obtained using the definite form and geometric-half identifications
from S77C1. The same-row identities transfer the grid list pin
at the original depth and bound. Keep the old witnesses $P,y,l$
of the remaining choice condition and compose their comparison
with that of the moved choice. This composition is componentwise
for the Galois, twisting and form data and transitivity of dependent
equality for the branch, twist, $\mathrm{sd},Q,k',\psi'$ entries.
This proves the character component on the specified moved data.

For the spectral component use the equality of the row packages,
transport the old spectral comparison for division change, and
identify the carriers of the resulting column forms with those
of the specified column. The extra condition from
$\mathrm{SExt.none}$ has its unique witness. For the object
component, the divisor in $O'$ is the transported one. Substitution
of the row identities in the division-geometry scalar conditions
transfers $x.ox$ to these rows. Together with the transported object
tie these four constructions prove $\mathrm{PositionMv}(p_0.s)$.

For the strengthened parameter $p_A$, the extension has the form
\[
  \mathrm{NEa}(h_L,\mathrm{NEo}(h_L,d),A).
\]
Its extra factor is exactly $A(K,B,f,T)$. Forget its witness
to apply the preceding construction. The division move keeps
$(K,B,f,T)$, so reinsert the given witness into the moved character
component. This includes the parameter with $\Gamma=\mathrm{anyGeometry}$
used in the common selection.
% END S77C2

% BEGIN S77C3
\paragraph{S77C3. A witness for the exact point represented by a good core.}
For either slot system $\mathcal S$ of S77C2, every context $K$ and
every $\beta:\mathrm{BPt}(K)$ satisfying
$\beta.\mathrm{Is}(\mathrm{IsGoodE}(\mathcal S))$, there exists
$q:\mathrm{WPt}(\mathcal S,K)$ such that
\[
 (\beta.D,\beta.O)=(q.w.D(q.Z),q.O),\qquad
             \mathrm{EntC}(q.x.\mathrm{ker}).
\]
The first equality belongs to the dependent sum $\sum_D\mathrm{Obj}(D)$.
The second says that the stored kernel datum $E$ equals the source
construction $\mathrm{kerEntC}$ for some place data, matrix family
and compatibility witness on this object.

First move a witness $x$ carrying this kernel property, for arbitrary
$N,h$. Transport its split-prime column and character data and use
them as the inputs of S77C2. Transfer the local division and theta--Weil
data along the same division change. The kernel family moves entry
by entry. Transfer its operators, its $E$, its standard data and
its local witnesses to that family. In a certificate
$E=\mathrm{kerEntC}(\mathrm{Pl},O_{\rm ker},B_s,H)$, copy the
place data and $H$, take the same matrix family $B_s$, and apply
the transport of $E$ to both sides of the equality. This proves
$\mathrm{EntC}$ on the resulting kernel. Apply the transport of
the entire tuple (group ties, comparison, twist, kernel family) to
$x.\mathrm{ker\_eq}$, giving the required kernel equality for
the moved object. Transport its error identity as well. The row
and four slot witnesses of S77C2, the transported local data,
and the division-change property of the good twists and theta
lattices now give a witness at $O.\mathrm{tr}(N,h)$.

Next unpack the definition of $\mathrm{IsGoodE}$ on $\beta$.
It presents a column, level, and core carrying a built record $b$,
$\mathrm{EntC}(b.\mathrm{ker})$ and the good-object condition.
The presentation identifies the \emph{whole} staged core with the
core of $\beta$. Substitute that equality before fixing the layer
$Y=\beta.Y$ and its size part $a=\beta.p$. The good-object condition
turns $b$ into the witness $b.\mathrm{wit}$ at its core object
$\mathrm{layObj}(Y)$. Apply the first part with
$N=\mathrm{nPrime}(a.L)$ and $h=a.\mathrm{ht}$. The point
$a.\mathrm{pt}$ has precisely this division change as its ambient
datum, and its object after forgetting the size presentation is
$\mathrm{layObj}(Y).\mathrm{tr}(N,h)$. Transport through these
two identities and take $q$ to consist of this base, frame, twist
datum, column, moved level, point object and moved witness. The
dependent-pair equality becomes reflexivity and the transported
kernel certificate proves the second assertion. This is the
source property $\mathrm{WitExE}(\mathcal S)$ for every $K$.
% END S77C3

% BEGIN S77C4
\paragraph{S77C4. Reciprocity data on the same point, before the bundle size.}
Assume the source hypotheses on discrete-series cover,
Harish--Chandra modules, Casselman--Wallach, Harish--Chandra
finiteness and cuspidal rapid decay, together with global Artin
reciprocity, the Weber ray-field statement and the source ray-field
hypothesis. Let $p$ be either parameter of S77C2 and take $\beta$
as in S77C3. There exist $q$ satisfying both conclusions of S77C3
and $R_q:\mathrm{RecT}(q.O.ex)$ with $\mathrm{RecPinE}(q,R_q)$.
The choice of $(q,R_q)$ can be made before any real bundle bound $c$.

Take $q$ from S77C3. The good-context certificate of $\beta$ makes
the level-prime set exactly
$\{g(0),g(1),p_{\rm de},p'_{\rm de},q_0\}$ in its host. The
core lies over the prescribed host and level. Its level formula is
the union of those primes and the image of the level choices. The
division change defining $\beta.D$ keeps the formula. This gives
$\mathrm{LevelTied}(\beta.D)$. Apply the first projection of
the dependent-pair equality in S77C3 to put it on $q.w.D(q.Z)$.

For zeta compatibility project the extension component successively
to $\mathrm{NEb}$: forget $A$ when present, then forget the extra
witness in $\mathrm{NEo}$ and the grid list in $\mathrm{NEl}$.
This supplies the original $\mathrm{BkLaw}$ on the same column.
In the source argument, use the kernel equality of $q.x$ to obtain
a tied kernel with the prescribed operators and kernel value. The
row extension then supplies the comparison columns. At the updated
face take its readings with the same tuple and ambient reading
$j_A$. These two equalities transfer the column evaluation in
$\mathrm{ReadC}$, keeping the word map and coordinate equivalences.
The division-pairing theorem on rational points for these exact
readings uses the five analytic hypotheses above. Apply the
$\mathrm{BkLaw}$ to that rationality and the reading-to-zeta
implication on the same row data. This gives
$\mathrm{ZetaCompatible}(q.w.D(q.Z))$.

Assemble the pinned linked-family data $S$ from the column,
character, Weil and twist data of $q$. For its reciprocity pieces
use the ray and modulus fields in the given host embedding,
as identified by Artin reciprocity and Weber's theorem, and use the
level and zeta conditions just proved. The kernel construction
on $q$, followed by its standard-reading refinement, gives a
kernel $k$ with the prescribed operators and $E$. S77C3 supplies
its entanglement identity. This produces $\mathrm{KPin}$ and both
standard-kernel assertions. Combine $S$ and $k$ with all coordinate
maps equal to the identity matrix. Their standard column, character
and twist comparisons verify the Gram, lattice, sharpness and frame
conditions on this record. The host pins from the ray-field
hypothesis and the original $q.x.\mathrm{pin}.\mathrm{d3b}$ now
supply host data for it. This gives $R_q$. Its two pins, its identity
coordinates, S77C3 and the two standard-kernel assertions are
exactly $\mathrm{RecPinE}(q,R_q)$.
% END S77C4

% BEGIN S77C5
\paragraph{S77C5. The remaining size-dependent step.}
For the same geometric parameter, assume the eventual bundle
implication on a prescribed witness and reciprocity datum of R73P5.
Then the witness and reciprocity inputs to the linked existence
step are supplied by S77C3--S77C4 at every good point. More
explicitly, restrict to the contexts for which the given eventual
implication holds for all $q,R_q,c$. Fix any prescribed good point
$\beta$ and use S77C4 before fixing $c$. Its valuation and scale
law give the required pair of nonnegative integers $(n_1,n_2)$ with
\[
 \beta.D.s.\mathrm{Del}=n_1,\qquad
 \beta.O.L.\mathrm{etaSz}=\beta.D.R.m\,n_2 .
\]
Transfer these identities, its valuation certificate and the scale
law along $(\beta.D,\beta.O)=(q.w.D(q.Z),q.O)$. Apply the
assumed implication for each real $c$ to this $q$ and this $R_q$.
It gives $b$ with $b.R=R_q$ and a witness $\mathrm{BWit}(q,b)$.
Transfer the entire triple $(q.w.D(q.Z),q.O,b)$ to the point of
$\beta$ using the same equality. In these coordinates the two
bundles compared by $\mathrm{linkOf}$ are the very same $b$.
Taking the large-base predicate to be identically true gives the
source $\mathrm{ExBL}$ implication. S77C4 is pointwise for every
context, so this step adds no eventual threshold to that of the
bundle implication.

The two larger existence requirements can now be located precisely.
S77C3 began with a good core, which already contains rows and their
readings. The family construction has to supply such cores at the
required points. In the final step the prescribed-$R_q$ bundle
implication uses the selected realization, geometric and spectral
comparisons, host and eigensystem inputs described in R73P5. B74R4
transports the rational tests and readings through a linked selection
once this existence and the family construction are supplied.
% END S77C5
